% from ega3-6-fr.tex
\oldpage[III]{165}
Pour simplifier l'écriture, nous nous bornerons au cas où $I$ est
l'intervalle $1\leq i\leq m$, et où la partition $(I_j)$ se compose des deux
intervalles $\{1,2,\ldots,r\}$ et $\{r+1,\ldots,m\}$.
\end{env}

\begin{proposition}[6.8.2]
Il existe un foncteur spectral canonique birégulier (dit « foncteur spectral
d'associativité ») noté
\[
  {}^{(e)}\mathscr E^S
  (f_1,\ldots,f_m;
   \mathscr P_\bullet^{(1)},\ldots,\mathscr P_\bullet^{(m)})
\]
(ou simplement ${}^{(e)}\mathscr E(f_1,\ldots,f_m;
\mathscr P_\bullet^{(1)},\ldots,\mathscr P_\bullet^{(m)})$) dont
l'aboutissement est
$\mathfrak{Tor}_\bullet^S(f_1,\ldots,f_m;
\mathscr P_\bullet^{(1)},\ldots,\mathscr P_\bullet^{(m)})$, et dont le terme
$E_2$ est donné par
\[
\begin{split}
  {}^{(e)}\mathscr E_{pq}^2
  =\bigoplus_{q_1+q_2=q}
  \mathscr{T}\!or_p^S\Bigl(
  &\mathfrak{Tor}_{q_1}^S
  (f_1,\ldots,f_r;
   \mathscr P_\bullet^{(1)},\ldots,\mathscr P_\bullet^{(r)}),\\
  &\mathfrak{Tor}_{q_2}^S
  (f_{r+1},\ldots,f_m;
   \mathscr P_\bullet^{(r+1)},\ldots,\mathscr P_\bullet^{(m)})
  \Bigr).
\end{split}
\]
\end{proposition}

\begin{proof}
Dans cet énoncé, on a identifié canoniquement $Y$ au produit
$Z^{(1)}\times_S Z^{(2)}$, où
$Z^{(1)}=Y^{(1)}\times_S Y^{(2)}\times_S\cdots\times_S Y^{(r)}$, et
$Z^{(2)}=Y^{(r+1)}\times_S\cdots\times_S Y^{(m)}$. Nous nous bornerons au
cas où $S$ et les $Y^{(i)}$ sont affines; on passe de ce cas particulier au
cas général par les méthodes développées dans (6.7.1) et (6.7.2), et nous
laissons les détails du raisonnement (sans difficulté) au lecteur. Nous
démontrerons donc le corollaire suivant.
\end{proof}

\begin{corollary}[6.8.3]
Soient $A$ un anneau, $S=\operatorname{Spec}(A)$, $X^{(i)}$
($1\leq i\leq m$) des $S$-schémas quasi-compacts et, pour chaque $i$, soit
$\mathscr P_\bullet^{(i)}$ un complexe de
$\mathscr O_{X^{(i)}}$-Modules quasi-cohérents limité inférieurement. Il
existe un foncteur spectral canonique birégulier ayant pour aboutissement
\[
  \mathbf{Tor}_\bullet^S
  (X^{(1)},\ldots,X^{(m)};
   \mathscr P_\bullet^{(1)},\ldots,\mathscr P_\bullet^{(m)})
\]
et dont le terme $E_2$ est donné par
\[
\begin{split}
  {}^{(e)}E_{pq}^2
  =\bigoplus_{q_1+q_2=q}
  \operatorname{Tor}_p^A\Bigl(
  &\mathbf{Tor}_{q_1}^S
  (X^{(1)},\ldots,X^{(r)};
   \mathscr P_\bullet^{(1)},\ldots,\mathscr P_\bullet^{(r)}),\\
  &\mathbf{Tor}_{q_2}^S
  (X^{(r+1)},\ldots,X^{(m)};
   \mathscr P_\bullet^{(r+1)},\ldots,\mathscr P_\bullet^{(m)})
  \Bigr).
\end{split}
\]
\end{corollary}

\begin{proof}
Suivant la définition donnée en (6.6.2), le calcul de l'hypertor considéré se
fait en prenant pour chaque $i$ un recouvrement ouvert affine fini
$\mathfrak U^{(i)}$ de $X^{(i)}$, en considérant les bicomplexes
$L_{\bullet\bullet}^{(i)}=C^\bullet(\mathfrak U^{(i)},
\mathscr P_\bullet^{(i)})$, une résolution projective
$M_{\bullet\bullet\bullet}^{(i)}$ de Cartan-Eilenberg de chacun de ces
bicomplexes (au sens de (0, 11.7.1)), le produit tensoriel
\[
  M_{\bullet\bullet\bullet}
  =\bigotimes_{i=1}^m M_{\bullet\bullet\bullet}^{(i)}
\]
de ces tricomplexes, et en prenant l'homologie de
$M_{\bullet\bullet\bullet}$. Considérons $M_{\bullet\bullet\bullet}$ comme
un complexe simple $N_\bullet$, produit tensoriel des deux complexes simples
\[
  N'_\bullet=\bigotimes_{i=1}^r M_{\bullet\bullet\bullet}^{(i)},
  \qquad
  N''_\bullet=\bigotimes_{i=r+1}^m M_{\bullet\bullet\bullet}^{(i)},
\]
où $N'_\bullet$ et $N''_\bullet$ sont gradués par la somme des degrés totaux
des $M_{\bullet\bullet\bullet}^{(i)}$. En outre, les $A$-modules des
complexes $N'_\bullet$ et $N''_\bullet$ sont projectifs, donc il résulte de
(6.5.9) que l'on a
\[
  H_\bullet(M_{\bullet\bullet\bullet})
  =\mathbf{Tor}_\bullet^A(N'_\bullet,N''_\bullet);
\]
la suite spectrale cherchée n'est autre alors que la suite (6.3.2.2)
appliquée aux complexes $N'_\bullet$ et $N''_\bullet$, compte tenu de
l'interprétation des modules d'homologie de ces complexes qui résulte de ce
qui précède (quand on applique les remarques du début à chacun des produits
partiels
$X^{(1)}\times_S\cdots\times_S X^{(r)}$ et
$X^{(r+1)}\times_S\cdots\times_S X^{(m)}$). Enfin, les propriétés de
régularité résultent de (6.3.2) et du fait que, les
$\mathscr P_\bullet^{(i)}$ étant limités inférieurement, il en est de même
des $M_{\bullet\bullet\bullet}^{(i)}$; par suite, $N'_\bullet$ et
$N''_\bullet$ sont limités inférieurement.
\end{proof}

